
qu'il résulte de 7.16.2 que @@M0@@ est isomorphe à un quotient d'un sous-groupe de @@M1@@ (@@M2@@ propre et lisse); en particulier, son rang est encore majoré par le rang de ce dernier.

\textbf{7.17.4.} Lorsque @@M3@@ est régulier en @@M4@@ et que @@M5@@ est maximal, on peut construire @@M6@@ en faisant éclater @@M7@@ dans le localisé de @@M8@@ en @@M9@@, et en localisant au point générique de la fibre de @@M10@@ dans l'éclaté. On trouve alors pour @@M11@@ une \emph{extension pure} de @@M12@@. Or il n'est pas difficile de montrer que si @@M13@@ est une extension pure d'un corps @@M14@@, alors pour tout schéma noethérien @@M15@@ sur @@M16@@, l'homomorphisme canonique @@M17@@ est un isomorphisme, en représentant @@M18@@ comme limite inductive des anneaux affines d'ouverts affines de schémas de la forme @@M19@@, et utilisant (IX 1.6). Donc on conclut de ceci que (7.17.1) peut être considéré comme un homomorphisme canonique
@@M20@@
Mais on notera qu'on ne peut en déduire un homomorphisme @@M21@@, car passant à une extension finie @@M22@@ de @@M23@@, et normalisant @@M24@@ dans @@M25@@, l'hypothèse de régularité sur @@M26@@ sera perdue.

Une autre façon de définir un homomorphisme (7.17.4.1) consiste à se ramener au cas où @@M27@@ est local de point fermé @@M28@@ (ce qui est trivial), et à considérer une suite croissante de sous-schémas fermés \emph{réguliers} @@M29@@ de @@M30@@, de point générique @@M31@@, tels que @@M32@@ (@@M33@@). Comme pour tout @@M34@@, @@M35@@, @@M36@@ est un anneau de valuation discrète de corps résiduel @@M37@@, corps des fractions @@M38@@, on dispose d'un homomorphisme de spécialisation
@@M39@@
